基于有向距离场（Signed Distance Field, SDF）的隐式几何表达在三维刚体碰撞检测中展现了极高的计算效率，但传统的一阶法向梯度接触模型在非平滑接触动力学（Non-Smooth Contact, NSC）中易因网格离散插值或切向大位移引起严重的局部速度突跃（毛刺）及约束发散。本文提出了一种基于函数空间高阶泰勒展开的二阶曲率感知平滑接触约束理论。通过提取体素 SDF 网格几何中的 Hessian 矩阵，将二阶曲率项以运动学预补偿参数的形式注入线性互补问题（LCP）方程的偏置项中。在此基础上，本文提出了拓扑感知流形降维分离（Topology-Aware Subsumption）、谱界限钳制（Spectral Bounding）及自适应权重衰减算法，将该框架无损侵入于现有的底层隐式求解器中。

通过高瞬态曲率的凸轮-从动件（Cam-follower）系统测试表明，本方法能够在纯刚体约束下完全平抑由离散分辨率导致的速度突跃伪影，在 $\sim 5$ 倍的超大积分时间步长（$\Delta t = 0.01\text{s}$）下依然体现出超越原生多边形网格模型（Native Mesh）的绝佳鲁棒性和 $\mathcal{O}(1)$ 量级的极速推断时间。同时，本文通过极限测试1:1双直齿轮刚性啮合（Twin Gear Benchmark），严格标定与评估了基于点切线导数构建物理空间场的一阶与二阶外推方法的数学极限。研究证明，在面对 $C^0$ 奇异非凸尖点与极小间隙的拓扑边界时，不可导面将诱发 Hessian 爆炸，导致任何基于局部微分算子的框架失效。基于此深刻的误差边界发现，本文在理论讨论的顶层蓝图中指出：未来面对高密复杂特征部件，接触动力学应逐步由点梯度范式（Point-based Differential）跃迁至体积积分SDF（Volumetric Integral SDF）范式，为真正的大型精细机械工业级虚拟组装验证建立了理论分水岭与演进路线。
